
\section{Theoretical Framework: Model Definition $m:n-CA^k$}

The flood simulation model is formalized using an extension of the Cellular Automata (CA) theory, known as $m:n-CA^k$ 
\parencite{Fonseca2011}. This architecture enables the integration of multiple layers of spatial and 
temporal information with different dimensions into a single cellular computational framework.

Furthermore, aligning with recent computational advancements, the model leverages tensorial data structures to provide a 
more accurate, flexible, and computationally efficient framework. This tensor-based approach allows for the 
modeling of continuous spatial domains and continuous cell states within the CA, simplifying the integration of complex 
geographical information \parencite{Fonseca2025}. 

Mathematically, the hydrodynamic cellular automaton is defined as a tuple:

$$A = \langle R, E, V, N, \Psi \rangle$$

Where:

\begin{itemize}
    \item $R$ is the $m$-dimensional cell space.
    \item $E$ is the set of $n$ layers that define the global state and the physical environment of the system.
    \item $V$ represents the vicinity function.
    \item $N$ represents the nucleus function.
    \item $\Psi$ is the combination or evolution function that governs the fluid dynamics.
\end{itemize}


\subsection{Cell Space ($R$) and Layer Definition ($E$)}

The simulation space $R$ is defined on a regular two-dimensional grid $R \subset \mathbb{Z}^2$, where each cell is 
identified by its spatial coordinates $(i, j)$. The physical size of each cell is determined by the spatial resolution 
$\Delta x$.

The system consists of a set $E = \{E_1, E_2, E_3, E_4\}$ of $n=4$ layers. One of the fundamental characteristics of the 
$m:n-CA^k$ model is that the layers may differ in their geometric and temporal dimensions. Through these tensor representations, 
they are defined as follows:

\begin{enumerate}
    \item $E_1$: Topobathymetry ($Z_{topo-bathy}$). A static 2D layer of dimension $X \times Y$. It combines the Digital 
    Elevation Model (DEM) with the bathymetry of the watercourses, providing the absolute elevation of the riverbed: 
    $E_1 \subset \mathbb{R}^{X \times Y}$.
    \item $E_2$: Land Cover (LC). A static 2D layer of size $X \times Y$. It identifies land uses, which are algebraically 
    mapped to Manning roughness coefficients ($\mu_{soil}$) via a state transfer function: 
    $E_2 \subset \mathbb{N}^{X \times Y} \xrightarrow{f} \mathbb{R}^{X \times Y}$.
    \item $E_3$: Water Depth ($WD$). A 2D dynamic layer of dimension $X \times Y$. It represents the main state of the cell 
    at time $t$, storing the water depth in meters: $E_3 \subset \mathbb{R}^{X \times Y}$.
    \item $E_4$: Rainfall ($RF$). A 3D dynamic layer of dimension $X \times Y \times T$, explicitly incorporating the time 
    domain to provide the rainfall rate for each step $\Delta t$: $E_4 \subset \mathbb{R}^{X \times Y \times T}$.
\end{enumerate}


\subsection{Vicinity Function ($V$) and Nucleus Function ($N$)}

The \textbf{Vicinity Function} ($V$) determines the subset of adjacent cells that influence the calculation of the future state of a 
given cell. A Moore neighborhood with radius $r=1$ is adopted, defining a set of 8 adjacent cells connected orthogonally and 
diagonally.

$$V(i,j) = \{ (i+u, j+v) \in R \mid \max(|u|, |v|) \le 1 \}$$

The \textbf{Nucleus Function} ($N$) is projected onto the central cell itself to update layer $E_3$, mapping the spatial 
modifications through local rules: $N(i,j) = \{ (i,j) \}$.


\subsection{Combination / Evolution Function ($\Psi$)}

The evolution function $\Psi: E^t \rightarrow E^{t+1}$ governs the transformation of the system from time $t$ to time $t+1$. 
In this proposed framework, the fluid dynamics computation is divided into two sequential phases: Meteorological Forcing and 
Hydrodynamic Routing.


\subsubsection{Phase 1: Meteorological Forcing}

At the start of iteration $t$, the water volume in the cell is updated by incorporating the contribution from the 3D rainfall 
layer ($E_4$). The intermediate state of the water depth, denoted $WD_{i,j}^*$, is calculated as:

$$WD_{i,j}^* = WD_{i,j}^t + RF(i,j,t)$$


\subsubsection{Phase 2: Kinematics, Viscosity, and Physical Limiters}

Instead of computing momentum conservation via dynamic forces, the evolution core utilizes a multi-physics Kinematic Wave 
approach combined with explicit spatial and temporal volume limiters. This formulation acts as a numerical safeguard against 
instabilities, ensuring accurate solutions when modeling shallow water flows over complex topographies, bottom discontinuities, 
and transient dry-bottom areas \parencite{Aleksyuk2017}.

For each cell in the neighborhood $k \in V(i,j)$, the water surface elevation ($WSE$) is evaluated as $WSE = E_1 + WD^*$. The 
hydraulic gradient $S_k$ towards each valid neighbor ($\Delta WSE_k > 0$) is calculated dividing the elevation drop by the 
directional distance $L_k$ ($\Delta x$ for orthogonal, $\Delta x\sqrt{2}$ for diagonal). The model identifies the maximum local 
slope $S_{max} = \max(S_k)$.

To accommodate fluids with varying rheological properties, the core evaluates the effective fluid velocity ($v_{approx}$) by 
continuously solving two competing flow regimes:

\begin{itemize}
    \item \textbf{Turbulent Regime:} Governs light, low-viscosity fluids (such as pure water). To this end, it is assumed that 
    the flow reaches its terminal velocity almost instantaneously. For the estimation of the kinematic velocity, the model adopts 
    the distributed uniform flow formulation proposed by \parencite{Bates2000} for raster-based flood inundation simulation. 
    Adapting the Manning's equation employed by the authors, and assuming the shallow water approximation where the hydraulic 
    radius is equivalent to the flow depth ($WD^*$), the velocity is defined as a function of the surface roughness coefficient 
    ($\mu_{soil}$) and the local slope gradient ($S_{max}$):
    $$v_{turb} = \frac{1}{\mu_{soil}} \cdot (WD^*)^{2/3} \cdot \sqrt{S_{max}}$$
    
    \item \textbf{Laminar/Viscous Regime:} Dictates the behavior of dense or non-Newtonian flows (such as mud or debris), driven 
    by fluid density ($\rho$) and dynamic viscosity ($\mu_{fluid}$):
    $$v_{lam} = \frac{\rho \cdot g \cdot (WD^*)^2 \cdot S_{max}}{3 \cdot \mu_{fluid}}$$
\end{itemize}

The automaton dictates that the fluid will travel at the lowest of the two velocities, naturally bounding the flow based on its 
dominant physical state:
$$v_{approx} = \min(v_{turb}, v_{lam})$$

Consequently, the theoretical outflow volume ($V_{theor}$) over the time step $\Delta t$ is calculated through the flow 
cross-section, and distributed into directional outflow volumes ($\Delta V_{theor \rightarrow k}$) applying Multiple Flow 
Direction (MFD) weights $W_k$:

$$V_{theor} = v_{approx} \cdot (WD^* \cdot \Delta x) \cdot \Delta t$$
$$\Delta V_{theor \rightarrow k} = V_{theor} \cdot \left( \frac{S_k}{\sum_{n \in V} S_n} \right)$$

To ensure strict mass conservation and prevent the inversion of hydraulic gradients (checkerboard instability), two sequential 
limiters are enforced before state updating:

\begin{enumerate}
    \item \textbf{Spatial Shield (Level Equalization):} The central cell is restricted from dropping its water level beyond half 
    the maximum elevation difference with its lowest neighbor:
    $$\Delta V_{spatial} = \frac{\max(\Delta WSE) \cdot \Delta x^2}{2}$$

    \item \textbf{Temporal Shield (Volumetric CFL Condition):} In traditional Computational Fluid Dynamics (CFD), the 
    Courant-Friedrichs-Lewy (CFL) condition is strictly enforced by constraining the time step ($\Delta t$) or limiting the 
    maximum velocity to ensure that a fluid wave does not traverse a distance greater than the grid resolution:
    $$v \le \frac{\Delta x}{\Delta t}$$
    
    Dynamically altering the global time step in parallel tensor executions is computationally inefficient. Instead, the 
    proposed model enforces the CFL condition volumetrically. By rearranging the classical kinematic condition 
    ($v \cdot \Delta t \le \Delta x$) and multiplying both sides by the flow cross-sectional area ($WD^* \cdot \Delta x$), 
    the inequality is transformed:
    $$(v \cdot \Delta t) \cdot (WD^* \cdot \Delta x) \le \Delta x \cdot (WD^* \cdot \Delta x)$$
    
    This mathematically maps to the model's volumetric routing variables, ensuring that the theoretical intended outflow 
    ($V_{theor}$) does not exceed the absolute total water volume present within the cell ($V_{cell} = WD^* \cdot \Delta x^2$). 
    Consequently, the effectively available volume ($V_{avail}$) is strictly bound by this physical fluid mass and the 
    aforementioned spatial restriction:
    $$V_{avail} = \min(WD^* \cdot \Delta x^2, \Delta V_{theor})$$
\end{enumerate}

A strict scale factor $\alpha$ is derived. If the intended routing volume exceeds the available constraints, all directional flows 
are scaled down proportionally:

$$\alpha = \frac{\min\left(\sum \Delta V_{theor \rightarrow k}, V_{avail}\right)}{\sum \Delta V_{theor \rightarrow k}}$$
$$\Delta V_{out \rightarrow k} = \Delta V_{theor \rightarrow k} \cdot \alpha$$

Finally, the main state variable ($E_3 \equiv WD$) is updated for $t+1$ by translating volumetric transfers ($m^3$) into vertical 
water column variations ($m$). A floating-point precision filter ($\epsilon = 10^{-5}$) is applied to avoid microscopic residual 
accumulations:

$$WD_{i,j}^{t+1} = \max\left(WD_{i,j}^* - \sum_{k \in V} \frac{\Delta V_{out \rightarrow k}}{\Delta x^2} + \sum_{k \in V} \frac{\Delta V_{in \leftarrow k}}{\Delta x^2}, \ 0\right)$$

This integrated tensorial update rigorously preserves volumetric continuity while seamlessly handling the wet-dry interface evolution.











\section{System Macro-Architecture}

The system is designed following a Loose Coupling strategy between two distinct technological domains, effectively 
separating "Data Engineering" (preparation) from "Computational Physics" (simulation). This decoupling is mediated 
by a standardized File-Based Interface, ensuring that the high-level data processing in Python and the high-performance 
execution in C++ remain independent yet perfectly synchronized.

\subsection{The Data-as-Interface Strategy}

Instead of integrating Python scripts directly into the C++ core via interpreters or APIs, the architecture utilizes 
a strictly defined File-Based Interface. The Data Pipeline acts as a "Producer", generating a structured repository 
of geographical and temporal data. The C++ Simulator acts as a "Consumer", ingesting these datasets through a dedicated 
InputAdapter.

The system is a unified environment based on the IDRISI standard. Static layers are stored as IDRISI Raster datasets 
(.doc/.img), while spatiotemporal data is managed through a custom Hierarchical IDRISI Format (HIF).

This strategy provides several key advantages:

\begin{enumerate}
    \item \textbf{Format Homogeneity}: By using a single underlying raster standard (IDRISI), the simulator's input 
    logic is simplified, reducing the overhead of managing multiple third-party libraries like HDF5 in the C++ 
    environment.
    \item \textbf{Technological Specialization}: Python manages complex GIS operations—such as coordinate reference 
    system (CRS) transformations and remote sensing API calls—while the C++ core focuses exclusively on the numerical 
    efficiency of Cellular Automata rules.
    \item \textbf{Immutability and Reproducibility}: Once the pipeline execution is finalized, the input artifacts 
    are fixed. This allows researchers to audit the exact state of the environment before a simulation run, ensuring 
    that results are deterministic and reproducible.
\end{enumerate}

\begin{figure}[htbp]
    \centering
    \includegraphics[width=0.85\textwidth]{chapter_3/system_architecture_overview}
    
    \caption{System architecture overview}
    \label{fig:system_architecture_overview}
\end{figure}


























